\documentclass[10pt]{mypackage}

\usepackage[uprightgreeks,noDcommand]{kpfonts}
\renewcommand{\coloneq}{\coloneqq}
\renewcommand{\eqcolon}{\eqqcolon}
%\usepackage{dsfont}
%\renewcommand{\mathbb}{\mathds}
%\usepackage{homework}
\usepackage{notes}

\usepackage[ backend=bibtex, style = alphabetic, sorting=ynt ]{biblatex}
\addbibresource{all_references.bib}

\usepackage{parskip}

\fancyhf{}
\fancyhead[R]{Avinash Iyer}
\fancyhead[L]{Gromov's Theorem on Groups of Polynomial Growth}
\fancyfoot[C]{\thepage}

\setcounter{secnumdepth}{0}

\begin{document}
\RaggedRight
This is going to be a broad overview of the ideas and techniques used to prove the following much celebrated theorem.
\begin{theorem}[Gromov's Theorem on Groups of Polynomial Growth]
  If a finitely generated group $\Gamma$ has polynomial growth, then $\Gamma$ contains a nilpotent subgroup of finite index.
\end{theorem}
All of this work is discussed in Gromov's paper, and I am only writing this up for personal use (and to present to the University of Virginia Graduate Students Seminar).
\section{Notions from Riemannian Geometry}%
\begin{definition}
  Let $f\colon X\rightarrow Y$ be a map. We call such a map \textit{globally expanding} if, for any two points $x_1\neq x_2$ in $X$, one has $d\left( f\left( x_1 \right),f\left( x_2 \right) \right) > d\left( x_1,x_2 \right)$.

  We say $f$ is expanding if each point $x\in X$ has a neighborhood $U$ such that $f|_{U}$ is globally expanding.
\end{definition}
Consider a compact connected Riemannian manifold, and let $f\colon X\to X$ be an expanding map. Then, $X$ has no boundary. There is then a map $ \widetilde{f}\colon \widetilde{X}\rightarrow \widetilde{X} $ induced on the universal cover. Since any globally expanding map is a covering map on itself, it follows that the map $ \widetilde{f} $ is a homeomorphism, and the inverse map $ \widetilde{f}^{-1}\colon \widetilde{X}\rightarrow \widetilde{X} $ is a contraction.

Furthermore, for any $\delta > 0$, there is some $\ve > 0$ such that for any $y_1,y_2\in \widetilde{X}$ with $d_{\widetilde{X}}\left( y_1,y_2 \right) \geq \delta$, one has
\begin{align*}
  d_{\widetilde{X}} \left( \widetilde{f}^{-1}\left( y_1 \right),\widetilde{f}^{-1}\left( y_2 \right) \right) &\leq \left( 1-\ve \right)d_{\widetilde{X}} \left( y_1,y_2 \right),
\end{align*}
which follows by the definition of globally expanding. Therefore, $ \widetilde{f}^{-1} $ has a unique fixed point by the contraction mapping theorem, so projecting back down yields that $f\colon X\rightarrow X$ has a fixed point as well.

Two theorems of Shub and Franks are fundamental to the proof. To understand them, though, we require understanding a few results.
\nocite{gromov_polynomial_growth_paper,van_den_dries_wilkie_paper}
\printbibliography
\end{document}
